Write \(\mathcal O_a=\{O\in\mathcal O:O\subseteq\mathcal U_a\}\) and \(m_a=|\mathcal O_a|\). Values of \(s_a\) outside \(\mathcal O_a\) are irrelevant. First suppose \(k\le M\). Given a score map \(s_a\), define \(O\preceq_a O'\) if and only if \(s_a(O)\le s_a(O')\). Then \[s_a(O_{T,a})\le Q_k(s_a(O_{1,a}),\ldots,s_a(O_{M,a}))\quad\Longleftrightarrow\quad\#\{i:s_a(O_{i,a})<s_a(O_{T,a})\}<k.\tag{1}\] Consequently the joint rank event depends only on the pair of total preorders induced by \(s_0,s_1\). Conversely, if a total preorder on \(\mathcal O_a\) has ordered equivalence classes \(C_{a,1}\prec_a\cdots\prec_a C_{a,d_a}\), assign score zero when \(d_a=1\), and otherwise assign \((j-1)K/(d_a-1)\) to every orbit in \(C_{a,j}\). Since \(K>0\), this is a representative in \([0,K]\) with \(d_a\le m_a\) levels. If \(k=M+1\), the event is identically true and the same finite reduction is immediate. The number of total preorders on an \(m\)-element labelled set is \[F_m=\sum_{d=1}^m d!\left\{\begin{matrix}m\\d\end{matrix}\right\},\tag{2}\] because one first partitions the set into \(d\) nonempty tie classes and then orders the classes. Thus only \(F_{m_0}F_{m_1}\) preorder pairs need be checked, proving finite attainment.

We now expose the exact objective and its complexity. Define \[\kappa_{w,a}(O\mid z)=\sum_{v:(v,a)\in O}\kappa_{w,a}(v\mid z),\qquad R_w(O_0,O_1)=\sum_z\nu(z)\kappa_{w,0}(O_0\mid z)\kappa_{w,1}(O_1\mid z).\tag{3}\] The product in (3) follows from conditional independence of the two auxiliary selector draws given the shared assignment. Independence across replicates makes the \(M\) calibration orbit pairs iid with mass \(R_w\). If \(O_a(v)\) denotes the orbit containing \((v,a)\), then, for a fixed preorder pair, its joint rank-event probability is exactly \[\sum_{v\in\mathcal V}q_{\mathrm{tar}}(v)\!\sum_{((O_{i,0},O_{i,1}))_{i=1}^M\in(\mathcal O_0\times\mathcal O_1)^M}\!\left\{\prod_{i=1}^M R_w(O_{i,0},O_{i,1})\right\}\prod_{a=0}^1\mathbf 1\!\left\{s_a(O_a(v))\le Q_k(s_a(O_{1,a}),\ldots,s_a(O_{M,a}))\right\}.\tag{4}\] Formula (4) retains the one coherent target member and the within-replicate cross-arm selector coupling; it assumes neither target-arm nor calibration-arm independence. Substitution of (4) into \(\texttt{def:coverage-frontier}\) gives the asserted evaluable finite maximum.

Let \(n=m_0m_1\). The \(n^M\) configuration weights in (4) can be tabulated by a depth-\(M\), \(n\)-ary recursion in \(O(n^M)\) arithmetic operations. For \(n\ge2\), enumerate configurations in reflected \(n\)-ary order, so successive configurations change one coordinate. For each target member and preorder pair, the two counts in (1) can then be updated in constant time after an \(O(M)\) initialization, and \(M=O(n^M)\). The case \(n=1\) is immediate. Multiplying by the numbers of preorder pairs and target members gives \[O\!\left(F_{m_0}F_{m_1}|\mathcal V|(m_0m_1)^M\right).\tag{5}\]

It remains to prove that binary scores need not attain the maximum. Take four pair blocks, \(\mathcal G=S_2^4\), and independent fair assignments with exactly one treated member per pair. Let \(\mathcal H\) be a singleton and \(w=1\). Given any assignment, each arm-\(a\) kernel chooses a block with probabilities \[q=\left(\frac{13}{20},\frac1{20},\frac14,\frac1{20}\right)\] and then selects that block's unique observed arm-\(a\) member. This kernel is observable and equivariant. Its block law is \(q\), independent of the assignment, and the auxiliary product rule gives \(R_w(b_0,b_1)=q_{b_0}q_{b_1}\). Set \[q_{\mathrm{tar}}(v_{b,r})=\frac{p_b}{2},\qquad p=\left(\frac1{20},\frac1{10},\frac14,\frac35\right).\] The two target lifts therefore share a block with law \(p\). Take \(M=3\) and \(\alpha=1/2\), so \(k=3\). Put \(s_0\equiv0\), and order the arm-\(1\) blocks as \(1\prec_1 3\prec_1 2\prec_1 4\), represented by \(0,K/3,2K/3,K\). Arm \(1\) fails exactly when all three calibration blocks are strictly below the target. Thus its failure probability is \[F_{\mathrm{multi}}=\frac14\left(\frac{13}{20}\right)^3+\frac1{10}\left(\frac{18}{20}\right)^3+\frac35\left(\frac{19}{20}\right)^3=\frac{104957}{160000}.\tag{6}\] Arm zero never fails, so the frontier objective is \(F_{\mathrm{multi}}-\alpha=24957/160000\).

For a binary map in arm \(a\), let \(H_a\subseteq\{1,2,3,4\}\) be the blocks assigned the high score \(K\), and put \(A_a=(1-q(H_a))^3\). Inclusion--exclusion conditional on the common target block gives the exact joint failure probability \[F_{\mathrm{bin}}(H_0,H_1)=p(H_0)A_0+p(H_1)A_1-p(H_0\cap H_1)A_0A_1.\tag{7}\] In the row order \(\varnothing,1,2,3,4,12,13,14,23,24,34,123,124,134,234,1234\), direct substitution of all sixteen possible \(H_1\) into (7) gives the following row maxima, over the common denominator \(320000000\): \[\begin{split}(&164616000,165302000,192052000,198366000,207906936,165912000,164712000,165787368,\\&203032000,208582976,208582976,164632000,165793875,164631423,202894331,164616000).\end{split}\tag{8}\] The table exhausts all \(16^2\) binary pairs. Hence \[\max_{H_0,H_1}F_{\mathrm{bin}}(H_0,H_1)=\frac{3259109}{5000000},\] attained, for example, by \(H_0=\{2,4\}\) and \(H_1=\{3,4\}\). The largest binary frontier objective is therefore \(759109/5000000\), and (6)--(8) yield the strict gap \[\frac{24957}{160000}-\frac{759109}{5000000}=\frac{83189}{20000000}>0.\] This proves that binary attainment is false.

Finally, suppose target and calibration orbit supports are disjoint in one arm and \(k\le M\). Assign score \(K\) to the target-supported orbits in that arm and zero elsewhere, and assign zero in the other arm. The adverse target score is always \(K\), every adverse calibration score is zero, and the joint rank event has probability zero. The frontier objective is at most \(1-\alpha\) by definition and attains that value here. Thus \(\Gamma_{M,\alpha}(w)=1-\alpha\), matching \(\texttt{prop:disjoint-orbit-witness}\). If \(k=M+1\), both rank thresholds in the frontier definition are \(+\infty\), the joint event has probability one for every score pair, and \(\Gamma_{M,\alpha}(w)=0\), matching \(\texttt{prop:rank-boundary}\).